% Generated draft: requires mathematical and source review.
\documentclass[11pt]{article}
\usepackage[margin=25mm]{geometry}
\usepackage{fontspec}
\setmainfont{Latin Modern Roman}
\usepackage{amsmath,amssymb,mathtools,unicode-math}
\setmathfont{Latin Modern Math}
\usepackage{enumitem,needspace,xurl,hyperref}
\setlength{\parindent}{0pt}
\setlength{\parskip}{5pt}
\setlength{\emergencystretch}{3em}
\newcommand{\N}{\mathbb{N}}
\newcommand{\Z}{\mathbb{Z}}
\newcommand{\Q}{\mathbb{Q}}
\newcommand{\R}{\mathbb{R}}
\newcommand{\C}{\mathbb{C}}
\newcommand{\F}{\mathbb{F}}
\newcommand{\A}{\mathbb{A}}
\newcommand{\PP}{\mathbb{P}}
\newcommand{\E}{\mathbb{E}}
\newcommand{\eps}{\varepsilon}
\newcommand{\dd}{\,\mathrm{d}}
\begin{document}
\section*{UnsolvedMath ID: 30004245}
\textbf{Source problem code:} catalogue entry
\subsection{Definitions and mathematical statement}
Consider the random lattice paths constructed from the height function of the two-periodic Aztec diamond, as referred to in the supplied source record. The record does not specify the tiling probability measure, periodic weights, boundary and height conventions, path construction, or ordering that defines a last path. Consequently, the following is a formulation template whose unresolved data must be supplied before it becomes a precise conjecture.

Let \(n\) denote the size of the diamond along admissible sizes, and let \(\theta_n\) denote its model parameters, possibly depending on \(n\). Once a path construction and ordering have been fixed, write \(\mathcal P_{n,\theta_n}\) for the resulting random path collection. A \emph{last path} means an extremal member of this collection with respect to that ordering; the ordering and the sense in which such a path is required to exist are not specified in the record. If this path admits a transverse-coordinate representation, denote that representation by \(L_{n,\theta_n}(u)\).

\textbf{Source question.} Can one establish the existence of a last path in this construction and prove that, after suitable centering and scaling, it converges to the Airy process? The supplied edited statement additionally asks for the parameter regimes in which this convergence holds.

Schematically, the requested convergence would have the form
\[
\left(
\frac{L_{n,\theta_n}(u_n+r_n t)-m_n(t)}{s_n}
\right)_{t\in I}
\ \Longrightarrow\
\bigl(\mathcal A(t)\bigr)_{t\in I},
\qquad n\to\infty,
\]
where \(u_n\) is an observation location, \(r_n>0\) and \(s_n>0\) are longitudinal and transverse scales, \(m_n\) is a centering function, \(I\subseteq\R\) is an observation interval, and \(\mathcal A\) denotes the intended Airy process. These are placeholders, not normalization data supplied by the source. A complete statement must specify the admissible parameter sequences \((\theta_n)\), observation region, deterministic or random nature of the normalization, interpolation or discrete evaluation convention, Airy-process variant, and meaning of convergence. In particular, the record does not determine whether convergence is meant only for finite-dimensional distributions or in a specified function-space topology.

The supplied dated assessment describes an original expectation covering all parameters, together with a later result for a last path in a thin rough--smooth regime. Neither the full parameter range nor the hypotheses and normalization of that restricted result are provided. No universal parameter quantifier or precise restricted theorem can therefore be extracted from this record alone.
\subsection{Short English statement}
The question asks whether the height-function paths of a two-periodic Aztec diamond have a well-defined last path whose centered and scaled fluctuations converge to an Airy process, and which parameter regimes permit this limit. The model, path ordering, normalization, and type of convergence need specification. The supplied status is partially solved, based on dated curation reporting a result in a thin rough--smooth regime; this draft does not independently verify that assessment.
\subsection{Sources}
\begin{itemize}
\item Local UnsolvedMath registry, permanent ID \(30004245\), source problem code \texttt{catalogue entry}, title \emph{Airy Limits for Last Paths in Aztec Diamonds}, in the supplied Oberwolfach Reports Open Problems collection. The registry records proposed year 2019 and supplies no individual proposer attribution.
\item Supplied source citation: \emph{Large Scale Stochastic Dynamics} (2020), Oberwolfach Reports, DOI \texttt{10.4171/owr/2019/42}, pp.\ 2605--2669. \url{https://doi.org/10.4171/owr/2019/42}. The background identifies Oberwolfach Reports as reports of mathematical workshops and their open-problem discussions. The supplied literature-source label for this reference is ``Oberwolfach report.''
\item Supplied literature source: \emph{Airy process at a thin rough region}. \url{https://arxiv.org/abs/2302.04663}.
\item Supplied literature source: \emph{Airy line ensemble at the boundary}. \url{https://arxiv.org/abs/2602.24063}.
\item Literature assessment checked at \texttt{2026-08-22}, as supplied dated curation, not independently verified: status ``partially solved.'' The supplied assessment states that later work proves Airy-process convergence for a last path in a thin rough--smooth regime, while the original all-parameter expectation is not fully established. The background explicitly requires independent verification of time-sensitive status and cited claims.
\item The registry labels the statement \texttt{corrected\_verified} and describes its editing as preserving mathematical content while repairing notation, typography, and sentence structure. This is a supplied editorial qualification, not independent mathematical verification of this draft or of the cited results.
\end{itemize}
\end{document}
